\documentclass[11pt,a4paper]{article}
\usepackage[utf8]{inputenc}
\usepackage[T2A]{fontenc}
\usepackage[russian,english]{babel}
\usepackage{amsmath,amssymb,booktabs,geometry}
\geometry{margin=25mm}
\title{Предсказательная точность усреднения и нелинейная устойчивость встречновращающихся RF-состояний}
\author{V17C.15 theory validation}
\date{21 ноября 2026 г.}
\begin{document}
\maketitle

\begin{abstract}
Проверяется, насколько leading усреднённые вращательные корни воспроизводятся прямой неавтономной динамикой идеального двухмодового RF-поля Local15. Введена безразмерная N-body постановка с разными массами, зарядами и трением; выведены условия коллективного вращательного баланса, Fourier-разложение внутренних мод регулярных колец и мера бассейна при заданном протоколе загрузки. Точная RF панель сохраняет оба знака вращения и предсказывает средние радиусы с ошибкой около 0.09\%, но частотная ошибка одного вида равна 8.386\%. Для одного заранее заданного старта RK4/128 согласуется с независимым DOP853 лучше 0.00089\% по радиусу и 0.00014\% по модулю частоты. Это проверка указанных моделей и протоколов, а не доказательство глобального аттрактора или высокой вероятности захвата.
\end{abstract}

\section{Безразмерная модель}
Положим \(\tau=\Omega_*t\), \(\mathbf x_i=\mathbf r_i/R\), \(\mu_i=m_i/m_*\), \(Q_i=q_i/q_*\), \(\lambda_i=\gamma_i/(m_i\Omega_*)\), и \(\kappa=q_*^2/(4\pi\varepsilon m_*R^3\Omega_*^2)\). Для (N) точечных частиц без внешних неэлектрических сил:
\begin{equation}
\mathbf x_i'=\mathbf u_i,\qquad
\mu_i\mathbf u_i'=-Q_i\nabla\phi(\mathbf x_i,\tau)-\mu_i\lambda_i\mathbf u_i
+\kappa\sum_{j\ne i}Q_iQ_j\frac{\mathbf x_i-\mathbf x_j}{|\mathbf x_i-\mathbf x_j|^3}.
\label{eq:nbody}
\end{equation}
Все частоты в безразмерной модели равны \(\nu_\ell=\omega_\ell/\Omega_*\). Ускорение (i) от парной силы включает делитель \(\mu_i\); зарядовая часть пропорциональна подписанному произведению \(Q_iQ_j\). Код `numerical/general_nbody.py` реализует конечный список signed-планарных гармоник и DC члены второго/четвёртого порядка; общая трёхмерная MFS/FEM передача пока остаётся интерфейсной задачей.

\section{Усреднение на движущейся орбите}
Для локально замороженной медленной координаты и одной гармонической силы
\[
\mathbf X''+\lambda\mathbf X'=\mathbf F_s(\mathbf X)+\operatorname{Re}\{\mathbf f(\mathbf X)e^{-i\nu\tau}\}
\]
быстрый отклик первого порядка равен
\[
\widehat{\boldsymbol\xi}=\frac{\widehat{\mathbf f}}{-\nu^2-i\lambda\nu},\qquad
\Delta\overline{\mathbf F}=\frac12\operatorname{Re}[(D\mathbf f)\widehat{\boldsymbol\xi}^{*}].
\]
Это асимптотика при малом микроперемещении, медленной огибающей и отсутствии близкого резонанса. При переходе \(\mathbf x=R_z(\Omega\tau)\mathbf y\) точная кинематика добавляет кориолисов, центробежный и вращательный drag члены:
\[
\mathbf y''+(\lambda I+2\Omega J_z)\mathbf y'
+(-\Omega^2P_{xy}+\lambda\Omega J_z)\mathbf y
=R_z(-\Omega\tau)\mathbf F(R_z(\Omega\tau)\mathbf y,\tau).
\]
Для гармоники с азимутальным индексом (k) частота равна \(\sigma=\nu-k\Omega\), однако одного изменения скалярного знаменателя недостаточно: ответ задаётся матрицей
\[
\mathcal D(\sigma)=-\sigma^2I-i\sigma(\lambda I+2\Omega J_z)
-\Omega^2P_{xy}+\lambda\Omega J_z-A_s.
\]
Формула движущейся орбиты выведена аналитически в `GENERAL_N_PARTICLE_THEORY_RU.md`, но количественное улучшение Local15 ею пока не установлено.

\section{Регулярные кольца и устойчивость}
Для кольца из (N_a) одинаковых частиц вида (s_a):
\[
\mathbf x_{aj}=R_a\mathbf e_r(\Theta_a+2\pi j/N_a)+Z_a\mathbf e_z.
\]
Суммарная кулоновская энергия одного кольца и внешняя сила на частицу:
\[
U_{aa}=\frac{\kappa Q_a^2N_a}{4R_a}\sum_{j=1}^{N_a-1}\csc\frac{\pi j}{N_a},\quad
F^{\mathrm{self}}_{R,a}=\frac{\kappa Q_a^2}{4R_a^2}\sum_{j=1}^{N_a-1}\csc\frac{\pi j}{N_a}.
\]
Для нескольких групп взаимодействие суммируется по всем (N_aN_b) парам; его производные задают радиальные/осевые силы и обменный момент, причём суммарный внутренний момент равен нулю. При циклической симметрии внутренние деформации раскладываются по дискретным Fourier-модам. Точная RF-устойчивость проверяется полной фундаментальной матрицей/мультипликаторами Флоке при периодическом режиме либо QR-показателями при несоизмеримых частотах. Эти спектральные свойства не определяют глобальную меру бассейна.

Для осевых смещений планарного одинаковозарядного кольца точные собственные значения кулоновского взвешенного Laplacian:
\[
\ell_m=\sum_{j=1}^{N-1}\frac{1-\cos(2\pi mj/N)}{[2R\sin(\pi j/N)]^3}.
\]
Линейная граница осевого buckling в квадрупольном DC-поле задаётся \((Q/\mu)2\beta_2=(\kappa Q^2/\mu)\max_m\ell_m\). Это локальный критерий для заданного планарного кольца; нелинейную конечную ветвь после потери жёсткости он не определяет. Точные суммы проверены прямым построением матриц для (N=2,3,6,8,12,16,24,32), максимум ошибки собственного спектра (1.45\times10^{-12}\).

\section{Прямая RF-проверка Local15}
Выполнено 48 интегрирований: 24 заранее зафиксированных старта на каждый вид, одна и та же позиционно-скоростная панель и RF-фазы, seed 15152026, RK4/96, (15\,000) периодов ведущей RF-моды. Модель содержит только идеальные (k=+2,-4) гармоники, линейное сопротивление и DC-квадруполь; Coulomb, реальные электроды и высшие MFS-моды отсутствуют. Потеря фиксируется при \(\rho\ge1\).

\begin{center}
\begin{tabular}{lrrrrrr}
\toprule
вид & \(\rho_\mathrm{lead}\) & \(\rho_\mathrm{RF}\) & ошибка \(\rho\) & \(\omega_\mathrm{lead}\) & \(\omega_\mathrm{RF}\) & ошибка \(|\omega|\)\\
\midrule
1 & 0.55036915 & 0.55085775 & +0.0888\% & +0.00144005 & +0.00156082 & +8.39\%\\
2 & 0.76638186 & 0.76569876 & -0.0891\% & -0.00105792 & -0.00105392 & -0.38\%\\
\bottomrule
\end{tabular}
\end{center}

По 22 из 24 стартов на каждый вид оставались удержанными; S08 и S16 теряются в первые (0.073) и (0.104) периода. Все 22 удержанные пары вращаются в противоположных направлениях. Фиксированная покрывающая панель не является iid, поэтому (22/24) не оценивает вероятность захвата.

Для S01 выполнено сравнение RK4/64, /96, /128 с DOP853 на 15 000 периодах и финальном 500-периодном окне. DOP853 применял \(\mathrm{rtol}=2\cdot10^{-11}\), \(\mathrm{atol}=2\cdot10^{-13}\) и максимальный шаг \(T_1/96\). RK4/128 даёт ошибки к DOP853: радиус −0.000405\% и частота +0.000883\% для вида 1; радиус +0.000121\% и частота −0.000139\% для вида 2. DOP853 для всей 24-стартовой панели не запускался.

\section{Бассейн и физическая параметризация}
Для структурного класса (S), меры начальных состояний и фаз \(\mu_\mathrm{load}\), и протокола загрузки \(\mathcal P_\mathrm{load}\), вероятность определяется как
\[
P_S(\mu_\mathrm{load})=\int\mathbf 1_{B_S}(X_0,\boldsymbol\psi_0;\mathcal P_\mathrm{load})\,d\mu_\mathrm{load}.
\]
Существование решения, локальная устойчивость и (P_S) различны. В сопротивлении Стокса с Cunningham mobility \(C_c\ge1\):
\[
\gamma=\frac{6\pi\eta a}{C_c},\qquad
\lambda=\frac{6\pi\eta a}{C_cm\Omega_*}.
\]
Архивная конвенция V17C.12B умножала \(6\pi\eta a\) на \(C_c\), поэтому old/new коэффициент drag равен \(C_c^2\). Парный Local15 тест с \(C_c=1.0027476353\), 24 фиксированными стартами на вид, RK4/96 и 3000 периодами удержал одни и те же 22/24 старта при обеих конвенциях; все удержанные знаки совпали. Это изолированная идеальная модель, не пересчёт физической V17C.12B/MFS панели. Исторические исходы остаются неизменными.

\section{Границы вывода}
Установлены алгебраические тождества общей модели, точные симметричные суммы и численная предсказательная ошибка конкретной Local15 панели. Не установлены полная неавтономная устойчивость, глобальные бассейны, MFS/FEM-гармонические поправки, универсальный порог коллективного перехода и физическая безопасность электродов. N=1,2,3,6 и другие конечные примеры являются проверками формул, не доказательством для произвольного (N).

\end{document}
